\section{Defining systems, initial states, and pulses}
\label{sec:systems}

\subsection{The system object}

\texttt{SpectroscopySystem} groups the Hamiltonian, detection dipole,
Lindblad collapse operators, and optional initial state. Its constructor is
\par\noindent
\begin{minipage}{\linewidth}
\begin{lstlisting}[style=python]
system = pulsus.SpectroscopySystem(
    H, dipole, collapse_ops=None, rho0=None, hbar=1.0,
    dipole_plus=None, dipole_minus=None,
)
\end{lstlisting}
\end{minipage}\par
The first two arguments are required. For Hilbert-space dimension $d$,
\texttt{H} and \texttt{dipole} must be dense $d\times d$ matrices. Supply the
same dimensions for each collapse operator, the initial density matrix, and
the two optional dipole components. NumPy arrays and nested Python lists
are suitable inputs; SciPy sparse matrices are not a supported input format
in this release. The implementation converts matrices to complex NumPy
arrays and constructs dense superoperators of size $d^2\times d^2$.

The field-free equation of motion is
\begin{equation}
 \dot\rho=-\frac{i}{\hbar}[H,\rho]
 +\sum_k\left(C_k\rho C_k^\dagger
 -\frac{1}{2}\{C_k^\dagger C_k,\rho\}\right).
\end{equation}
The interaction superoperator uses
$\mathcal V\rho=(i/\hbar)[\mu,\rho]$. The sign and Fourier conventions are
therefore already built into PULSUS; do not insert additional factors of
$i$ into a spectrum simply to match an assumed plotting convention.

The object exposes \texttt{H}, \texttt{dipole}, \texttt{collapse\_ops},
\texttt{rho0}, \texttt{hbar}, and \texttt{d}. Its precomputed
\texttt{L} and \texttt{V} are the field-free and full interaction
superoperators. When dipole components are supplied, \texttt{V\_plus} and
\texttt{V\_minus} are constructed as well. Rebuild the object when changing
the Hamiltonian, dipole, dissipation, or $\hbar$: editing an attribute does
not automatically recompute the corresponding superoperators.

\subsection{Collapse operators contain their rates}

For a decay channel with rate $\gamma$ and lowering operator $a$, supply
$C=\sqrt{\gamma}\,a$, rather than $\gamma a$. The dissipator is quadratic in
$C$. The dimer tutorial also includes excitation with occupation parameter
$\bar n$:
\par\noindent
\begin{minipage}{\linewidth}
\begin{lstlisting}[style=python]
gamma = 0.05
nbar = 0.05
collapse_ops = [
    np.sqrt(gamma * (nbar + 1.0)) * a,
    np.sqrt(gamma * nbar) * a.conj().T,
    np.sqrt(gamma * (nbar + 1.0)) * b,
    np.sqrt(gamma * nbar) * b.conj().T,
]
\end{lstlisting}
\end{minipage}\par
Here $a$ and $b$ act on their respective sites in the full tensor-product
space. A single-site $2\times2$ collapse operator cannot be inserted into
the four-state dimer without the tensor-product extension. These local
excitation and decay terms define the model's bath; their stationary state
should not be assumed to equal a Gibbs state of an arbitrary coupled
Hamiltonian.

\subsection{Choose one consistent unit system}

PULSUS does not convert units. All arguments must be consistent with the
equations above. In the tutorial, $\hbar=1$: Hamiltonian entries,
\texttt{omega\_L}, \texttt{omega1}, \texttt{omega3}, \texttt{eta}, and
decay rates use inverse model-time units; \texttt{sigma}, \texttt{T}, and
time axes use the corresponding model-time unit. Frequencies are angular
frequencies, not cycles per unit time.

For energies in eV and times in fs, one consistent choice is
\par\noindent
\begin{minipage}{\linewidth}
\begin{lstlisting}[style=python]
hbar_eV_fs = 0.6582119569
system = pulsus.SpectroscopySystem(
    H=H_eV,
    dipole=dipole,
    collapse_ops=collapse_ops_per_sqrt_fs,
    hbar=hbar_eV_fs,
)
omega_L = photon_energy_eV / hbar_eV_fs
omega_axis = energy_axis_eV / hbar_eV_fs
\end{lstlisting}
\end{minipage}\par
In this convention, \texttt{sigma} and \texttt{T} are in fs and
\texttt{eta} and $\gamma$ are in fs$^{-1}$. A population lifetime
$\tau$ gives $\gamma=1/\tau$ in the chosen time unit. Include
$\sqrt\gamma$ in each collapse operator once. The code divides Hamiltonian
and interaction commutators by $\hbar$; it does not divide the Lindblad
dissipator by $\hbar$. Do not apply an extra $1/\hbar$ to rates that already
have inverse-time units. The product of dipole and field amplitude must
also have energy units if absolute physical polarization amplitudes are
required. The illustrative examples use model dipoles and fields.

\subsection{The state before the first pulse}

\texttt{rho0} is the system's density matrix immediately before the first
interaction. Its diagonal entries give populations; its off-diagonal
entries describe coherences. For a $d$-state model, supply a complex
$d\times d$ matrix with trace one.

The tutorials start from a stationary state: a state that does not change
under the Hamiltonian and dissipation before the pulses arrive. Calculate
that state once, then pass it to the spectrum calculation:
\par\noindent
\begin{minipage}{\linewidth}
\begin{lstlisting}[style=python]
rho_ss = system.stationary_state()
S = pulsus.finite_pulse_spectrum(
    system=system,
    pulse1=pulse, pulse2=pulse, pulse3=pulse,
    omega1=omega1, omega3=omega3, T=20.0, eta=0.02,
    pathway="NR", rho0=rho_ss,
)
\end{lstlisting}
\end{minipage}\par

The method returns the matrix; it does not assign
\texttt{system.rho0}. An explicit \texttt{rho0} in a function call takes
priority over a state stored in the system. If neither is supplied,
PULSUS raises an error rather than choosing a state for you.

The standard finite-pulse formulation assumes stationary preparation.
For a model with several possible stationary states, choose the state
that matches the intended preparation. The solver returns one stationary
candidate and does not determine whether it is unique.

\subsection{Dipole components for the rotating-wave approximation}

Full-interaction calculations require only \texttt{dipole}. RWA
calculations additionally require both \texttt{dipole\_plus} and
\texttt{dipole\_minus}. PULSUS does not infer them from \texttt{H}.
For the excitation-conserving dimer in the first tutorial,
\par\noindent
\begin{minipage}{\linewidth}
\begin{lstlisting}[style=python]
dipole_plus = a.conj().T + b.conj().T
dipole_minus = dipole_plus.conj().T
dipole = dipole_plus + dipole_minus
system = pulsus.SpectroscopySystem(
    H=H, dipole=dipole, collapse_ops=collapse_ops,
    hbar=1.0, dipole_plus=dipole_plus,
    dipole_minus=dipole_minus,
)
rho_ss = system.stationary_state()
\end{lstlisting}
\end{minipage}\par
The raising/lowering definition must match the physical transition
structure of the user's model. A matrix that happens to be triangular in
an arbitrary basis is not automatically its positive-frequency dipole.
The RWA functions use \texttt{V\_plus} for field sign $+1$ and
\texttt{V\_minus} for sign $-1$; detection still uses the full dipole.

\subsection{GaussianPulse and pulse duration}

\par\noindent
\begin{minipage}{\linewidth}
\begin{lstlisting}[style=python]
pulse = pulsus.GaussianPulse(
    omega_L=1.10, sigma=1.25, E0=1.0, phase=0.0,
)
\end{lstlisting}
\end{minipage}\par

\texttt{omega\_L} is the carrier angular frequency and \texttt{sigma}
is the temporal width of the field envelope. Both are required.
\texttt{E0} is the peak field amplitude (default 1.0), and
\texttt{phase} is in radians (default 0.0). To change pulse parameters,
create a new \texttt{GaussianPulse}.

PULSUS uses the field envelope $\exp[-t^2/(2\sigma^2)]$.
If a pulse duration is given as a full width at half maximum (FWHM),
convert it to \texttt{sigma} before creating the pulse:
\begin{equation}
 \sigma=\frac{\mathrm{FWHM}_{E}}{2\sqrt{2\ln2}}
       \simeq\frac{\mathrm{FWHM}_{E}}{2.355},
 \qquad
 \sigma=\frac{\mathrm{FWHM}_{I}}{2\sqrt{\ln2}}
       \simeq\frac{\mathrm{FWHM}_{I}}{1.665}.
\end{equation}
Here $E$ means the \emph{field-amplitude envelope} and $I$ means the
\emph{intensity envelope}, proportional to the squared field envelope.
Use the formula matching the stated duration; neither FWHM is
\texttt{sigma} itself. Larger \texttt{sigma} gives a longer pulse with a
narrower frequency bandwidth.

When changing the duration, keeping \texttt{E0} fixed keeps the peak
field fixed. The pulse area then changes with \texttt{sigma}. To compare
fixed envelope areas instead, choose $E_0\propto1/\sigma$. Finite-pulse
and short-pulse results include field-amplitude factors; impulsive
molecular responses do not. Normalizing each plot separately removes
those amplitude differences from the display.
